% Draft (2026-11-07): the evidence ceiling. To be revised after GPT review round P2.
\section{The ceiling is a property of the evidence}
\label{sec:framework}

\paragraph{Setting.}
An audit case consists of a program $X$, the evidence $E$ shown to the auditor (the task text, the visible tests,
any added text, and anything the auditor can compute from these, including by executing $X$), and a hidden
contract $H$ that defines the ground truth. The program is defective, $D=1$, if it violates $H$. An auditor is any
measurable, possibly randomised, map from $E$ to a flag; we write $\phi(e)=\Pr(\text{flag}\mid E=e)\in[0,1]$.
Sampling more readings, using a larger model, reasoning longer or executing the supplied code changes $\phi$ but not
the joint law of $(E,D)$. Write $P_1$ and $P_0$ for the laws of $E$ given $D=1$ and $D=0$.

\begin{proposition}[Evidence ceiling]
\label{prop:ceiling}
For every auditor $\phi$,
\[
\mathrm{TPR}(\phi)-\mathrm{FPR}(\phi)\;\le\;\mathrm{TV}(P_1,P_0),
\]
and for every $\alpha\in[0,1]$, $\max\{\mathrm{TPR}(\phi):\mathrm{FPR}(\phi)\le\alpha\}$ is attained by a
likelihood-ratio test of $E$, so the receiver operating characteristic of the likelihood ratio is the ceiling for
all auditors that see $E$.
\end{proposition}

\begin{proof}
$\mathrm{TPR}-\mathrm{FPR}=\int\phi\,(p_1-p_0)\,d\mu\le\int_{p_1>p_0}(p_1-p_0)\,d\mu=\mathrm{TV}(P_1,P_0)$ for
densities $p_1,p_0$ with respect to a dominating measure $\mu$, since $0\le\phi\le1$. The second statement is the
Neyman--Pearson lemma.
\end{proof}

\begin{corollary}[Indistinguishable component]
\label{cor:component}
Suppose the cases split into components and, on a component $K$, $E$ has the same law for defective and correct
programs. Then every auditor has $\mathrm{TPR}_K=\mathrm{FPR}_K$. If a share $\pi$ of the defects lies in such
components, recall at component false-flag rate at most $\alpha$ is at most $(1-\pi)+\pi\alpha$.
\end{corollary}

Both statements are textbook facts about hypothesis testing (the Neyman--Pearson lemma and the total-variation
bound on any test's power minus its size; \citealp{neyman1933efficient}); we claim no novelty for them. Note that the
component bound needs the false-flag constraint on the component itself: under an overall false-flag rate $\alpha$
alone, with $q$ the share of correct programs in the component, the bound is only $(1-\pi)+\pi\min(1,\alpha/q)$.
Their use here is to separate two causes of a missed defect that audit evaluations conflate. A miss on a distinguishable case is a \emph{capacity} miss: some auditor with the same evidence could do
better. A miss on an indistinguishable case is an \emph{information} miss: no auditor can do better without
different evidence. Only interventions that change $E$ move the ceiling.

\paragraph{Contract randomisation makes the component observable.}
On a natural benchmark the hidden contract is fixed to the benchmark's reference behaviour, so indistinguishability
is a counterfactual: we cannot see the auditor's response under the other admissible contract. We construct it.
Take a discrepancy unit: an input class on which two behaviours, $A$ and $B$, are both admissible under the task
text and the visible tests. Build two programs that differ only on that class, $X_A$ and $X_B$, and draw
$H\in\{A,B\}$ independently of everything the auditor sees.

\begin{lemma}
\label{lem:symmetric}
If $E$ depends on $(X,H)$ only through $X$ and $H$ is independent of $X$ with $\Pr(H=A)=\tfrac12$, then $P_1=P_0$ on
the unit, so every auditor's recall equals its false-flag rate, whatever its capacity.
\end{lemma}

\begin{proof}
With $X$ uniform over $\{X_A,X_B\}$, $D=1$ exactly on $(X_A,B)$ and $(X_B,A)$ and $D=0$ on $(X_A,A)$ and $(X_B,B)$.
Each event has probability $\tfrac14$, and the law of $E$ given each depends only on $X$, so $P_1=P_0=\tfrac12
Q_{X_A}+\tfrac12 Q_{X_B}$, where $Q_x$ is the law of $E$ given $X=x$.
\end{proof}

Revealing $H$ in the evidence (stating the requirement) breaks the premise of Lemma~\ref{lem:symmetric}; text that
draws attention to the input class, or lists both behaviours without choosing, does not. The experiment in
Section~\ref{sec:e6} therefore measures three things that the propositions leave open: how far real auditors sit
below the ceiling once the requirement is revealed, whether additional computation helps more before or after
the reveal, and how much of the recall measured on natural benchmarks is agreement with the benchmark's convention
on indistinguishable units, which would be false flags had the convention gone the other way.

\paragraph{What we do not claim.}
The ceiling is relative to the evidence, and the evidence includes what a model brings with it: a model that has
memorised a benchmark's reference sees more than the task text. Contract randomisation neutralises this on the
constructed units, because the memorised convention is equally likely to be wrong. On natural units we can only
bound the information-limited share using human judgements of admissibility, which we report with their agreement.
